\subsection{部分积}

设 \((\mathbf{X}_i)_{i\in I}\) 是一族集合，并设 \(J\) 为 \(I\) 的子集。积 \(\prod_{i\in J} X_i\) 称为 \(\prod_{i\in I} X_i\) 的部分积。如果 \(f\) 是一个函数，其图像 \(F\) 是 \(\prod_{i\in I} X_i\) 的成员，那么 \(F \circ \Delta_J\) （其中 \(\Delta_J\) 是 \(J \times J\) 的对角线 ) 是 \(f\) 到 \(J\) 的限制的图。显然 \(F \circ \Delta_J \in \prod_{i\in J} X_i\) ；映射 \(F \to F \circ \Delta_J\) （从 \(\prod_{i\in I} X_i\) 到 \(\prod_{i\in I} X_i\) ）称为索引为 \(J\) 的投影，并记作 \(\operatorname{pr}_J\) .

命题5. 设 \((\mathbf{X}_{\iota})_{\iota \in \mathbf{I}}\) 是一族集合，且设 \(\mathbf{J}\) 为 \(\mathbf{I}\) 。若对每个 \(\iota \in \mathbf{I}\) 我们拥有 \(\mathbf{X}_{\iota} \neq \emptyset\) ，投影 \(\operatorname{pr}_{\mathbf{J}}\) 是……的映射 \(\prod_{\iota \in \mathbf{I}} \mathbf{X}_{\iota}\) 到上 \(\prod_{\iota \in \mathbf{J}} \mathbf{X}_{\iota}\) .

鉴于上述评论，只需证明以下命题即可：

命题6. 设 \((\mathbf{X}_{\iota})_{\iota \in \mathbf{I}}\) 为一个集合族，使得 \(\mathbf{X}_{\iota} \neq \emptyset\) 对于所有 \(\iota \in \mathbf{I}\) 成立。如果 \(g\) 是从 \(\mathbf{J} \subset \mathbf{I}\) 到 \(\mathbf{A} = \bigcup_{\iota \in \mathbf{I}} \mathbf{X}_{\iota}\) 的映射，使得 \(g(\iota) \in \mathbf{X}_{\iota}\) 对于所有 \(\iota \in \mathbf{J}\) 成立，那么存在一个 \(f\) ，它是 \(g\) 到 \(\mathbf{I}\) 的扩张，使得 \(f(\iota) \in \mathbf{X}_{\iota}\) 对于所有 \(\iota \in \mathbf{I}\) 成立.

对于每一个 \(\iota\in I-J\) 让 \(T_{\iota}\) 表示项 \(\tau_{y}(y\in X_{\iota})\) 。由于 \(X_{\iota}\neq\emptyset\) 根据假设，我们有 \(T_{\iota}\in X_{\iota}\) 对于所有 \(\iota\in I-J\) （第一章，第4节，第1号）。若G是 \(g\) 的图形，则图形 \(G\cup\left(\bigcup_{\iota\in I-J}\{(\iota,T_{\iota})\}\right)\) 是一个具有所需性质的函数的图形，这一点可以立即验证。

推论1. 设 \((\mathbf{X}_{\iota})_{\iota \in \mathbf{I}}\) 设为一族集合，使得对每个 \(\iota \in \mathbf{I}\) 我们拥有 \(\mathbf{X}_{\iota} \neq \emptyset\) 然后对每一个 \(\alpha \in \mathbf{I}\) 投影 \(\operatorname{pr}_{\alpha}\) 是……的映射 \(\prod_{\iota \in \mathbf{I}} \mathbf{X}_{\iota}\) 到上 \(\mathbf{X}_{\alpha}\) .

将命题5应用于子集 \(J = \{\alpha\}\) 的 I 并注意到 \(\mathbf{pr}_{\alpha}\) 是以下映射的复合： \(X_{\alpha}^{\{\alpha\}}\) 到上 \(X_{\alpha}\) 与映射的复合 \(\mathbf{pr}_{\{\alpha\}}\) .

推论2。设 \((\mathbf{X}_t)_{i \in \mathbf{I}}\) 是一个集合族。则 \(\prod_{i \in \mathbf{I}} \mathbf{X}_t = \varnothing\) 当且仅当存在 \(i \in \mathbf{I}\) 使得 \(\mathbf{X}_t = \varnothing\) .

如果我们有 \(X_{t} \neq \varnothing\) 对于每一个 \(t \in I\) ，则由推论1可知 \(\prod_{i \in I} X_{t} \neq \varnothing\) ；反之，如果 \(\prod_{i \in I} X_{t} \neq \varnothing\) ，关系 \(\Pr_{\alpha}\left(\prod_{i \in I} X_{t}\right) \subset X_{\alpha}\) 表明 \(X_{\alpha} \neq \varnothing\) 对于每一个 \(\alpha \in I\) .

因此，如果我们有一个族 \((\mathbf{X}_{t})_{t\in\mathbf{I}}\) 对于非空集合族，我们可以引入（作为辅助常项）一个定义域为I的函数f，使得 \(f(t)\in\mathbf{X}_{t}\) 对于所有 \(t\in I\) 在实践中人们常说：“取一个元素 \(x_{t}\) 在每个 \(X_{t}\) 直观地说，我们因此“选择”了一个元素。 \(x_{t}\) 在每个集合中 \(X_{t}\) ；逻辑符号的引入 \(\tau\) 并且规范其使用的准则使我们无需制定一个“选择公理”来使这一操作合法化（参见结果概要，第4节，第10条）。

推论3。设 \((\mathbf{X}_t)_{i\in I}\) 并且 \((\mathbf{Y}_t)_{i\in I}\) 是两族具有相同指标集I的集合。如果 \(\mathbf{X}_t\subset \mathbf{Y}_t\) 对于每一个 \(i\in I\) ，那么

\[
\prod_ {i \in I} X _ {i} \subset \prod_ {i \in I} Y _ {i}.
\]

反之，如果 \(\prod_{i\in I}X_i\subset \prod_{i\in I}Y_i\) ，并且如果 \(X_{i}\neq \emptyset\) 对于每一个 \(i\in I\) ，那么 \(X_{i}\subset Y_{i}\) 对于每一个 \(i\in I\) .

第一个断言是显然的，第二个断言则来自命题1的推论5，因为此时对每个 \(\alpha \in I\) ,

\[
\mathrm{X} _ {\alpha} = \operatorname{pr} _ {\alpha} \left(\prod_ {i \in \mathrm{I}} \mathrm{X} _ {i}\right) \subset \operatorname{pr} _ {\alpha} \left(\prod_ {i \in \mathrm{I}} \mathrm{Y} _ {i}\right) = \mathrm{Y} _ {\alpha}.
\]

